% Created 2026-01-26 Mon 20:10
% Intended LaTeX compiler: pdflatex
\documentclass[11pt]{article}
\usepackage[utf8]{inputenc}
\usepackage[T1]{fontenc}
\usepackage{graphicx}
\usepackage{longtable}
\usepackage{wrapfig}
\usepackage{rotating}
\usepackage[normalem]{ulem}
\usepackage{amsmath}
\usepackage{amssymb}
\usepackage{capt-of}
\usepackage{hyperref}
\author{John Zimmerman}
\date{2026:01:22}
\title{Demo - Web Foundations}
\hypersetup{
 pdfauthor={John Zimmerman},
 pdftitle={Demo - Web Foundations},
 pdfkeywords={},
 pdfsubject={},
 pdfcreator={Emacs 30.2 (Org mode 9.7.11)}, 
 pdflang={English}}
\begin{document}

\maketitle
\tableofcontents

\section{Demo - Web Foundations}
\label{sec:org3c4655f}

\subsection{The Beginning}
\label{sec:orge9c3088}

\textbf{Tim Berners-Lee}, a british scientist at CERN, invented to World Wide Web in 1989. Back then, the web was conceived and developed to meet the demand for automatic information-sharing between scientists in universities and institutions.
\subsection{How the Internet Works}
\label{sec:orgf5e9aa7}

Very simply, On the client side there is a brower like Chrome of Firefox and it pushes requests to the server, and the server sends a response back to the browser. Somewhere down the line, the browser creates the \textbf{Document Object Model} (DOM) and the data (HTML/CSS) shows up visually on the screen.
\subsection{The Trinity of Web Development}
\label{sec:org9ed8df5}

The three main components of web developnment is \textbf{HTML}, \textbf{CSS}, and \textbf{JavaScript}.

\begin{description}
\item[{HTML}] HyperText Markup Language
\begin{itemize}
\item Builds the \emph{structure} or skeleton of website
\item Think of a frame on a house
\end{itemize}
\item[{CSS}] Cascadimg Style Sheets
\begin{itemize}
\item Builds the \emph{presentation} of a website
\item Think of the decorations on a house
\end{itemize}
\item[{JS}] JavaScript
\begin{itemize}
\item Builds the \emph{behavior} and interactivity of a website
\item Like the appliances of a house
\end{itemize}
\end{description}
\subsection{The Basic Structure of a Webpage}
\label{sec:org91d9f79}

\begin{verbatim}
<!DOCTYPE html> 
<html> 
<head> 
  <title>Jen's Kitchen</title> 
</head> 

<body> 
<h1>This is a heading</h1>

<h2> this is another heading</h2>

<p>This is a paragraph</p>

<br> 

</body> 
</html> 
\end{verbatim}
\subsection{HTML Structure}
\label{sec:org4411280}

\begin{description}
\item[{Tag}] Refers to the descriptor or type for content. Example <p>, <h1>, etc.
\item[{Content}] Refers to the \emph{stuff} in between the angle brackets of a tag.
\item[{Attribute}] Goes inside of a tag, used to add clarity and nuance to an element.
\item[{Value}] Content for the attribute.
\item[{Element}] The tag and the content together is the element.
\end{description}

If a tag doesn't have content, it's considered a \textbf{self-closing tag}. All other tags need a \emph{closing tag}. Example. </p>, </h1>, etc.

\begin{verbatim}
// examples here...
\end{verbatim}
\subsection{WWW}
\label{sec:orga604708}

For websites, we need a website server. This means software that handles HTTP or HTTPS. Servers job is to \emph{serve up} documents.
\subsection{IP Addresses}
\label{sec:org5add84e}

\begin{verbatim}

- Protocol :: https://
- Host name :: www.
- Domain Name :: Name of the site
- Dir path :: /directories/
- Doc :: absolute path - index.html

https://www.nameofsite/directories/index.html

\end{verbatim}
\subsection{Default files}
\label{sec:orga6b3ec5}

The index.html is usually the file in the root of the directory and is the main document for the webpage. If you don't specify a path, it will default to index.html.
\end{document}
